\documentclass[10pt,a4paper]{amsart}
\usepackage{amsmath,amssymb,amsthm,graphicx}
\usepackage{verbatim}
\newenvironment{reference}{\comment}{\endcomment}
%\newenvironment{reference}{}{}
\newenvironment{slogan}{\comment}{\endcomment}
\newenvironment{history}{\comment}{\endcomment}

% For commutative diagrams we use Xy-pic
\usepackage[all]{xy}

% We use 2cell for 2-commutative diagrams.
\xyoption{2cell}
% Theorem environments.
%
\theoremstyle{plain}
\newtheorem{theorem}[subsection]{Theorem}
\newtheorem{proposition}[subsection]{Proposition}
\newtheorem{lemma}[subsection]{Lemma}

\theoremstyle{definition}
\newtheorem{definition}[subsection]{Definition}
\newtheorem{example}[subsection]{Example}
\newtheorem{exercise}[subsection]{Exercise}
\newtheorem{situation}[subsection]{Situation}

\theoremstyle{remark}
\newtheorem{remark}[subsection]{Remark}
\newtheorem{remarks}[subsection]{Remarks}

\numberwithin{equation}{subsection}

% Macros
%
\def\lim{\mathop{\mathrm{lim}}\nolimits}
\def\colim{\mathop{\mathrm{colim}}\nolimits}
\def\Spec{\mathop{\mathrm{Spec}}}
\def\Hom{\mathop{\mathrm{Hom}}\nolimits}
\def\Ext{\mathop{\mathrm{Ext}}\nolimits}
\def\SheafHom{\mathop{\mathcal{H}\!\mathit{om}}\nolimits}
\def\SheafExt{\mathop{\mathcal{E}\!\mathit{xt}}\nolimits}
\def\Sch{\mathit{Sch}}
\def\Mor{\mathop{\mathrm{Mor}}\nolimits}
\def\Ob{\mathop{\mathrm{Ob}}\nolimits}
\def\Sh{\mathop{\mathit{Sh}}\nolimits}
\def\NL{\mathop{N\!L}\nolimits}
\def\CH{\mathop{\mathrm{CH}}\nolimits}
\def\proetale{{pro\text{-}\acute{e}tale}}
\def\etale{{\acute{e}tale}}
\def\QCoh{\mathit{QCoh}}
\def\Ker{\mathop{\mathrm{Ker}}}
\def\Im{\mathop{\mathrm{Im}}}
\def\Coker{\mathop{\mathrm{Coker}}}
\def\Coim{\mathop{\mathrm{Coim}}}

% Boxtimes
%
\DeclareMathSymbol{\boxtimes}{\mathbin}{AMSa}{"02}

%
% Macros for moduli stacks/spaces
%
\def\QCohstack{\mathcal{QC}\!\mathit{oh}}
\def\Cohstack{\mathcal{C}\!\mathit{oh}}
\def\Spacesstack{\mathcal{S}\!\mathit{paces}}
\def\Quotfunctor{\mathrm{Quot}}
\def\Hilbfunctor{\mathrm{Hilb}}
\def\Curvesstack{\mathcal{C}\!\mathit{urves}}
\def\Polarizedstack{\mathcal{P}\!\mathit{olarized}}
\def\Complexesstack{\mathcal{C}\!\mathit{omplexes}}
% \Pic is the operator that assigns to X its picard group, usage \Pic(X)
% \Picardstack_{X/B} denotes the Picard stack of X over B
% \Picardfunctor_{X/B} denotes the Picard functor of X over B
\def\Pic{\mathop{\mathrm{Pic}}\nolimits}
\def\Picardstack{\mathcal{P}\!\mathit{ic}}
\def\Picardfunctor{\mathrm{Pic}}
\def\Deformationcategory{\mathcal{D}\!\mathit{ef}}
\usepackage{cleveref}
\crefname{originallemma}{original lemma}{original lemmas}
\Crefname{originallemma}{Original lemma}{Original lemmas}
\makeatletter
\newlabel{limits-lemma-descend-syntomic}{{0C3L}{source}{}{sourcepointer-1}{}}
\newlabel{limits-lemma-descend-syntomic@cref}{{[originallemma][0][]0C3L}{[1][0][]source}{}{}{}}
\makeatother
\crefname{originallemma}{original lemma}{original lemmas}
\Crefname{originallemma}{Original lemma}{Original lemmas}
\makeatletter
\newlabel{limits-lemma-descend-unramified}{{0C4W}{source}{}{sourcepointer-2}{}}
\newlabel{limits-lemma-descend-unramified@cref}{{[originallemma][0][]0C4W}{[1][0][]source}{}{}{}}
\makeatother
\crefname{originallemma}{original lemma}{original lemmas}
\Crefname{originallemma}{Original lemma}{Original lemmas}
\makeatletter
\newlabel{limits-lemma-limit-contained-in-constructible}{{05F4}{source}{}{sourcepointer-3}{}}
\newlabel{limits-lemma-limit-contained-in-constructible@cref}{{[originallemma][0][]05F4}{[1][0][]source}{}{}{}}
\makeatother
\crefname{originallemma}{original lemma}{original lemmas}
\Crefname{originallemma}{Original lemma}{Original lemmas}
\makeatletter
\newlabel{morphisms-lemma-dimension-fibre-after-base-change}{{02FY}{source}{}{sourcepointer-4}{}}
\newlabel{morphisms-lemma-dimension-fibre-after-base-change@cref}{{[originallemma][0][]02FY}{[1][0][]source}{}{}{}}
\makeatother
\crefname{originallemma}{original lemma}{original lemmas}
\Crefname{originallemma}{Original lemma}{Original lemmas}
\makeatletter
\newlabel{morphisms-lemma-syntomic-relative-dimension}{{02K1}{source}{}{sourcepointer-5}{}}
\newlabel{morphisms-lemma-syntomic-relative-dimension@cref}{{[originallemma][0][]02K1}{[1][0][]source}{}{}{}}
\makeatother
\title{At-worst-nodal relative dimension 1 descends along directed limits of affine transition morphisms}
\author{The Stacks Project Authors}
\date{stacks-project:curves.tex:local-snapshot-20261012; GNU FDL 1.2 or later}
\begin{document}
\maketitle
\noindent\textit{Editorial note.} All delivered mathematical text is the author's own, copied from the pinned Stacks Project chapter file curves.tex; no statement or proof step was rewritten, reordered or removed.\par\medskip
\noindent\textit{Editorial note.} The counted proof is the author's complete argument: the syntomic descent, the relative-dimension decomposition, the constructible containment descent, and the unramifiedness descent are all retained verbatim.\par\medskip
\noindent\textit{Editorial note.} The label lemma-descend-nodal-family names the counted statement, and the proof environment immediately following it in the source is delivered verbatim as its proof.\par\medskip
\noindent\textit{Editorial note.} The counted statement is carried by the author's own lemma environment in the chapter source.\par\medskip
\noindent\textit{Allowed standard prerequisites (editorial).}\begin{itemize}
\item Descent of the syntomic property. Limits, Lemma limits-lemma-descend-syntomic. Produces an index i at which the base-changed morphism f\_i is syntomic.
\item Syntomic relative dimension and constructible containment. Morphisms, Lemmas morphisms-lemma-syntomic-relative-dimension and morphisms-lemma-dimension-fibre-after-base-change; Limits, Lemma limits-lemma-limit-contained-in-constructible. Decomposes X\_i into open-and-closed relative-dimension pieces, shows the limit X maps into the dimension-1 piece, and descends that containment to some finite level.
\item Descent of unramifiedness. Limits, Lemma limits-lemma-descend-unramified. Promotes the unramifiedness of Sing(f) over Y, known after base change to Y, back to some finite level i'.
\end{itemize}\medskip
\ifdefined\hypertarget\hypertarget{sourcepointer-1}{}\fi
\noindent\textit{Source reference.} Original lemma 0C3L: sources/stacks-chow/tags.txt, source lines 15196--15196. Prerequisite pointer invoked at proof line(s) 5113; it names a standard fact of a sibling chapter that the counted argument uses, and every step of the counted proof built on it is retained verbatim in the delivered body.\par\medskip
\ifdefined\hypertarget\hypertarget{sourcepointer-2}{}\fi
\noindent\textit{Source reference.} Original lemma 0C4W: sources/stacks-chow/tags.txt, source lines 15241--15241. Prerequisite pointer invoked at proof line(s) 5130; it names a standard fact of a sibling chapter that the counted argument uses, and every step of the counted proof built on it is retained verbatim in the delivered body.\par\medskip
\ifdefined\hypertarget\hypertarget{sourcepointer-3}{}\fi
\noindent\textit{Source reference.} Original lemma 05F4: sources/stacks-chow/tags.txt, source lines 7013--7013. Prerequisite pointer invoked at proof line(s) 5124; it names a standard fact of a sibling chapter that the counted argument uses, and every step of the counted proof built on it is retained verbatim in the delivered body.\par\medskip
\ifdefined\hypertarget\hypertarget{sourcepointer-4}{}\fi
\noindent\textit{Source reference.} Original lemma 02FY: sources/stacks-chow/tags.txt, source lines 3209--3209. Prerequisite pointer invoked at proof line(s) 5120; it names a standard fact of a sibling chapter that the counted argument uses, and every step of the counted proof built on it is retained verbatim in the delivered body.\par\medskip
\ifdefined\hypertarget\hypertarget{sourcepointer-5}{}\fi
\noindent\textit{Source reference.} Original lemma 02K1: sources/stacks-chow/tags.txt, source lines 3356--3356. Prerequisite pointer invoked at proof line(s) 5118; it names a standard fact of a sibling chapter that the counted argument uses, and every step of the counted proof built on it is retained verbatim in the delivered body.\par\medskip
\begin{lemma}
\label{lemma-descend-nodal-family}
Let $S = \lim S_i$ be a limit of a directed system of schemes
with affine transition morphisms.
Let $0 \in I$ and let $f_0 : X_0 \to Y_0$ be a morphism of schemes over $S_0$.
Assume $S_0$, $X_0$, $Y_0$ are quasi-compact and quasi-separated.
Let $f_i : X_i \to Y_i$ be the base change of $f_0$ to $S_i$ and
let $f : X \to Y$ be the base change of $f_0$ to $S$.
If
\begin{enumerate}
\item $f$ is at-worst-nodal of relative dimension $1$, and
\item $f_0$ is locally of finite presentation,
\end{enumerate}
then there exists an $i \geq 0$ such that $f_i$ is at-worst-nodal
of relative dimension $1$.
\end{lemma}

\begin{proof}
By Limits, Lemma \ref{limits-lemma-descend-syntomic}
there exists an $i$ such that $f_i$ is syntomic.
Then $X_i = \coprod_{d \geq 0} X_{i, d}$ is a disjoint union of
open and closed subschemes such that $X_{i, d} \to Y_i$
has relative dimension $d$, see
Morphisms, Lemma \ref{morphisms-lemma-syntomic-relative-dimension}.
Because of the behaviour of dimensions of fibres under base change given in
Morphisms, Lemma \ref{morphisms-lemma-dimension-fibre-after-base-change}
we see that $X \to X_i$ maps into $X_{i, 1}$.
Then there exists an $i' \geq i$ such that $X_{i'} \to X_i$
maps into $X_{i, 1}$, see
Limits, Lemma \ref{limits-lemma-limit-contained-in-constructible}.
Thus $f_{i'} : X_{i'} \to Y_{i'}$ is syntomic of relative dimension $1$
(by Morphisms, Lemma \ref{morphisms-lemma-dimension-fibre-after-base-change}
again).
Consider the morphism $\text{Sing}(f_{i'}) \to Y_{i'}$.
We know that the base change to $Y$ is an unramified morphism.
Hence by Limits, Lemma \ref{limits-lemma-descend-unramified}
we see that after increasing $i'$ the morphism
$\text{Sing}(f_{i'}) \to Y_{i'}$ becomes unramified.
This finishes the proof.
\end{proof}

\end{document}
